\section{Introdução}

Em sistemas embarcados, o uso da tecnologia de reconhecimento de palavras-chave (\textit{keyword spotting}) é amplamente utilizado nos dias atuais, permitindo que o usuário acione dispositivos sem o uso das mãos. As expressões mais populares são ``OK Google'', ``Alexa'' e ``Hey Siri''. Assim, esses dispositivos ficam em escuta contínua até que a palavra-chave seja identificada, para que possam se conectar à internet e executar alguma ação solicitada pelo usuário~\cite{carlos2025}.

Para que esses dispositivos possam reconhecer essas palavras, é necessário embarcar algoritmos de Aprendizagem de Máquina (AM). Todavia, há o desafio de inserir esses modelos nesses \textit{hardwares} (microcontroladores), devido aos seus recursos limitados de memória e processamento. Por conta disso, surge o \textit{Tiny Machine Learning} (TinyML), focado em reduzir e comprimir os modelos de AM~\cite{carlos2025,pereira2023}.

Existem várias técnicas a fim de reduzir os modelos de AM, entre elas, destaca-se a quantização, que mapeia os pesos, vieses e ativações de ponto flutuante (FP32) para representações inteiras de baixa precisão~\cite{ravaglia2021,zhao2021}. Para quantizar um modelo, há duas abordagens principais: o \textit{Post-Training Quantization} (PTQ) e o \textit{Quantization-Aware Training} (QAT). O primeiro quantiza os parâmetros após o treinamento do modelo, utilizando estratégias de calibração como \textit{Min-Max} ou \textit{Mean Squared Error} (MSE)~\cite{zhao2021}. Já no segundo, durante o treinamento, os parâmetros são quantizados de forma simulada (\textit{fake quantization}), permitindo que a rede se ajuste por meio do \textit{Straight-Through Estimator} (STE) considerando o erro de discretização e a acurácia~\cite{bhalgat2020,lin2023}. Embora a quantização para 8~bits (compressão de $4\times$ em relação ao FP32) já apresente resultados satisfatórios com PTQ~\cite{carlos2025,ravaglia2021}, a redução para precisões \textit{sub-8-bit} --- como 4~bits ($8\times$) e 2~bits ($16\times$) --- tende a degradar severamente o desempenho sem reajuste dos parâmetros~\cite{bhalgat2020,lin2023}. Nesse cenário, variantes do QAT baseadas em gradiente, como o \textit{Learned Step Size Quantization} (LSQ) e o LSQ+, formulam a quantização como um problema de otimização conjunta, aprendendo também o passo de quantização e o deslocamento (\textit{offset}) durante o treinamento, diferentemente de abordagens que mantêm esses parâmetros fixos~\cite{bhalgat2020}.

Assim, o presente artigo tem como objetivo investigar essas duas abordagens de quantização tomando como base o trabalho ``Sistema Embarcado para Reconhecimento de Palavras usando ESP32-S3 e TinyML'', de Carlos~\cite{carlos2025}. Nesse estudo, o autor utilizou a plataforma Edge Impulse desde o pré-processamento dos dados até a implantação no microcontrolador, aplicando apenas a técnica PTQ para converter os parâmetros de ponto flutuante para inteiros de 8~bits. Diante disso, em vez de utilizar o Edge Impulse, este trabalho reproduz em PyTorch o mesmo \textit{pipeline} de extração de características MFCC e a arquitetura 1D-CNN (1.974 parâmetros), mantendo idênticas as parametrizações originais para focar exclusivamente na etapa de quantização. Com isso, busca-se investigar se há vantagem em utilizar o QAT em sua forma padrão e em suas variantes LSQ, LSQ+ e LSQ+ Híbrido frente ao PTQ (\textit{Min-Max} e \textit{MSE}). Para tanto, pesos e ativações são quantizados para 8~bits, 4~bits e 2~bits, comparando-se os resultados por meio das métricas de acurácia, função de perda (\textit{loss}), AUC e \textit{F1-score} ponderado.

O restante deste artigo está organizado da seguinte forma: a Seção~II apresenta a fundamentação teórica sobre quantização e formulação de otimização das técnicas PTQ e QAT avaliadas; a Seção~III descreve a metodologia, o conjunto de dados e a configuração experimental em PyTorch; a Seção~IV analisa os resultados comparativos nas precisões de 8, 4 e 2~bits; por fim, a Seção~V apresenta as conclusões e trabalhos futuros.
